\documentclass[10pt,a4paper]{amsart}
\usepackage[margin=19mm]{geometry}
\usepackage{amsmath,amssymb,mathtools}
\usepackage[scr=rsfs,cal=euler]{mathalfa}
\usepackage{xfrac}
\usepackage[unicode,hidelinks,bookmarks=false]{hyperref}
\newcommand{\ie}{i.e.\ }
\newcommand{\cf}{cf.\ }
\newcommand{\resp}{respectively\ }
\newcommand{\Subsection}{\subsection}
\providecommand{\nobreakdash}{\nobreak\hbox{-}}
\setlength{\emergencystretch}{2em}
\multlinegap0pt
\usepackage{bm}
\usepackage{bbm}
\usepackage{dsfont}
\usepackage{xspace}
\usepackage{cancel}
\usepackage{mathtools}
\usepackage{amsthm}
\makeatletter
\@ifundefined{theorem}{\newtheorem{theorem}{Theorem}}{}
\@ifundefined{lemma}{\newtheorem{lemma}[theorem]{Lemma}}{}
\makeatother
\title[Combinatorial Structure of Inert and Ambiguous Classes in Mo]{Combinatorial Structure of Inert and Ambiguous Classes in Modular Group}
\author{Debattam Das, Krishnendu Gongopadhyay,, Khushi Mishra}
\date{Source: arXiv:2602.19176v1 (CC BY 4.0)}
\begin{document}
\maketitle

\noindent\emph{Source and licence.} Combinatorial Structure of Inert and Ambiguous Classes in Modular Group. Version arXiv:2602.19176v1 (CC BY 4.0), distributed under the Creative Commons Attribution 4.0 International licence, as recorded in the arXiv licence metadata for this exact version and shown on the abstract page https://arxiv.org/abs/2602.19176v1. Full rights notice: licenses/arxiv-2602.19176-CC-BY-4.0.txt.\par\smallskip

\noindent\textit{Editorial scope.} Counted unit: the Lemma whose source label is \texttt{2.1} in the article (source0.tex lines 334--356, source label \texttt{2.1}), delivered together with the article's own complete original proof of it (source0.tex lines 358--456, one \texttt{proof} environment of 99 lines). No mathematics is written, repaired or re-derived: statement and proof are the article's own text line for line. Required context retained uncounted: none. Every definition, display and label of those lines is retained unchanged. Omitted from this package and not redistributed: 1--333, 357--357, 457--1123. Nothing inside a retained excerpt is omitted or altered: every retained body is delivered exactly as the article's lines are written, so no omission receipt of any kind accompanies this package. Citations: the retained excerpts carry no citation command at all, so no bibliography entry of the article is transcribed and no reference is redistributed. Reference adaptation: none was needed and none was made; every reference inside the delivered unit resolves inside it, so no unresolved reference or citation is produced. Printed slips kept literal: no printed word, symbol, hypothesis or number of the counted statement or of its proof was altered. Layout and class: the article's own class is \texttt{amsart} with packages including geometry, amsmath, amsfonts, amssymb, xcolor, mathrsfs, etoolbox, array, hyperref, caption; the wrapper is the lane's pinned \texttt{amsart} editorial wrapper at 10pt on A4, which restates the theorem-like environments the retained text needs and adds the article's own macro definitions that the retained text needs (none), with extra wrapper packages (bm, bbm, dsfont, xspace, cancel, mathtools). The delivered numbering is the unit's own. Assets: no retained span contains a figure environment, a graphics-inclusion command, a PDF-page inclusion, a table or a TikZ picture; no image file is copied into this package, the package declares no asset dependency and no image file is redistributed with it. Licence: version arXiv:2602.19176v1 of this article is distributed under the Creative Commons Attribution 4.0 International licence, as recorded in the arXiv licence metadata for this exact version and shown on the abstract page https://arxiv.org/abs/2602.19176v1. A full rights notice is delivered at licenses/arxiv-2602.19176-CC-BY-4.0.txt. Characters: every retained excerpt is pure ASCII, so the delivered engine, XeLaTeX with its default Latin Modern text font, reports no missing character.
\begin{lemma}\label{2.1}
Let $\{p\}$ be an inert class in $\mathscr{I}_{4t}$, and suppose
$\{p\}$ contains a cyclically reduced word of the form
\[
p \;=\; AB^{\alpha_1}AB^{\alpha_2}\cdots AB^{\alpha_{2t}}, 
\qquad \alpha_i \in \{\pm 1\}.
\]
If $\{p\}$ is fixed by $\phi_w$, then there exists an integer $m$ with 
$1 \le m \le t$ and $m \mid t$ such that, after a cyclic rotation of $p$, we have
\[
p \;=\;
\big(AB^{\alpha_1}\cdots AB^{\alpha_m}\big)
\big(AB^{-\alpha_1}\cdots AB^{-\alpha_m}\big)
\big(AB^{\alpha_1}\cdots AB^{\alpha_m}\big)
\big(AB^{-\alpha_1}\cdots AB^{-\alpha_m}\big)\cdots\cdots,
\]
that is,
\[
\alpha_{m+i} = -\alpha_i \qquad \text{for all $i$ (indices modulo $2t$)}.
\]
Equivalently, the $2t$--tuple $(\alpha_1,\dots,\alpha_{2t})$ has period $2m$
and the word $p$ is a $(t/m)$--fold repetition of a word of length $2m$.
\end{lemma}

\begin{proof}
Suppose $\{p\}\in \mathscr I_{4t}$ and let
$$
p \;\longleftrightarrow\;
\begin{pmatrix}
a & b\\
c & d
\end{pmatrix},
\qquad ad-bc=1,\quad |a+d|>2.
$$
Then
$$
\phi_w(p)=w^{-1}pw
\;\longleftrightarrow\;
\begin{pmatrix}
a & -b\\
-c & d
\end{pmatrix}.
$$

Moreover,
$$
\phi_w(A)
\;\longleftrightarrow\;
\phi_w\!\left(
\begin{pmatrix}
0 & -1\\
1 & 0
\end{pmatrix}\right)
=
\begin{pmatrix}
0 & 1\\
-1 & 0
\end{pmatrix}
\;\longleftrightarrow\; A,
$$
and
$$
\phi_w(B)
\;\longleftrightarrow\;
\phi_w\!\left(
\begin{pmatrix}
0 & -1\\
1 & 1
\end{pmatrix}\right)
=
\begin{pmatrix}
0 & 1\\
-1 & 1
\end{pmatrix}
\;\longleftrightarrow\; AB^{-1}A.
$$
Consequently,
$$
\phi_w(A)=A
\qquad\text{and}\qquad
\phi_w(B)=AB^{-1}A.
$$
Here $p \in\langle~ P ~\rangle$, since neither $A$ nor $B$ is hyperbolic, p cannot be conjugate to either of them. Moreover, every $\gamma \in W$ that is not conjugate to a generator is conjugate to an $(AB)$-word. Therefore, every conjugacy class of $\mathscr{I}_{4t}$ contains an AB-word.\\ Let $p$ itself be an $(AB)$-word of length 4t, that is $p$ = $AB^{\alpha_1}AB^{\alpha_2}....AB^{\alpha_{2t}}$, where $\alpha_i$'s are $\pm1$. $\phi_w$ is an inner automorphism, therefore $$\phi_w(AB^{\alpha_1}AB^{\alpha_2}\dots AB^{\alpha_{2t}}) = B^{-\alpha_1}AB^{-\alpha_2}\dots AB^{-\alpha_{2t}}A.$$
 Since $\{p\}\in \mathscr{I}_{4t}$, the conjugacy class $\{p\}$ is fixed by the automorphism 
$\phi_{w}$. This means that the words $p$ and $\varphi_{w}(p)$ must be conjugate in the 
group. The word $p$ has the standard form
\[
p = AB^{\alpha_1}\, AB^{\alpha_2}\cdots AB^{\alpha_{2t}},
\]
and applying $\phi_{w}$ sends each block $AB^{\alpha_i}$ to $B^{-\alpha_i}A$. Hence 
$\phi_{w}(p)$ begins with $B$ ends with $A$ and has the structure of a $(BA)$-word, while $p$ begins 
with $A$ ends with $B$ and is an $(AB)$-word. A $(BA)$-word can be conjugate to an $(AB)$-word only after 
an odd cyclic rotation: even rotations preserve the initial $B$, whereas odd rotations 
convert the leading $BA$ into $AB$. Therefore, any conjugating element must implement an odd 
cyclic shift of $\varphi_{w}(p)$.

Consider $L_m = B^{\alpha_m}A\, B^{\alpha_{m-1}}A \cdots B^{\alpha_1}A$. Conjugation by $L_m$ performs exactly the odd cyclic permutation 
required to convert the $(BA)$-pattern of $\varphi_{w}(p)$ back into the $(AB)$-pattern of 
$p$, giving the identity
\[
L_m\,\varphi_{w}(p)\,L_m^{-1} = p,
\]
which shows that $\varphi_{w}(p)$ and $p$ represent the same conjugacy class.

Moreover, writing out the equality $L_m\,\varphi_{w}(p)\,L_m^{-1} = p$ block-by-block yields a 
constraint on the exponents $\alpha_i$, i.e.,
$$ 
(B^{\alpha_m}A....B^{\alpha_2}AB^{\alpha_1})(B^{-\alpha_1}AB^{-\alpha_2}....AB^{-\alpha_{2t}}A)(B^{-\alpha_1}AB^{-\alpha_2}....AB^{-\alpha_m}) = AB^{\alpha_1}AB^{\alpha_2}....AB^{\alpha_{2t}}.
 $$   
 Then, $$\alpha_{m+1} = -\alpha_1,
 ~\alpha_{m+2} = -\alpha_2,
 \ldots, ~\alpha_{m+m} = -\alpha_m,~
 \alpha_{m+(m+1)} = -\alpha_{m+1},~
 \ldots, \alpha_{m} = -\alpha_t.$$
Matching corresponding $(AB)$-blocks on both sides 
forces the sign in the $(m+i)$-th position to be the negative of the sign in the $i$-th 
position. Hence
\[
\alpha_{m+i} = -\alpha_i \qquad \text{for all } i \text{ (indices taken modulo $2t$). }
\]
In particular, the sequence $(\alpha_1,\ldots,\alpha_{2t})$ becomes a repetition of a block of 
length $2m$, so the total length $2t$ must be an integer multiple of $2m$. Therefore, $m \mid t$.
\end{proof}

\end{document}
